\section{Modelo de datos}

Todo vive en un único proyecto Supabase (\cd{jijklguopafevyucogro}). El modelo tiene
cuatro capas: el crudo se guarda intacto y cada capa siguiente es una vista que agrega
información sin destruir la anterior.

\begin{figure}[htbp]
\centering
\resizebox{0.98\textwidth}{!}{% Figura: capas del modelo de datos. El crudo nunca se toca; cada capa agrega.
\begin{tikzpicture}[node distance=7mm]

\node[almacen,text width=46mm] (crudo) {\textbf{Capa 0 --- crudo}\\
  \cd{monitoreo\_sc\_electrico}\\ \cd{radiacion\_sc\_15s}\\
  \nota{tal como salió del CSV}};

\node[comp,text width=46mm,right=14mm of crudo] (corr) {\textbf{Capa 1 --- corrección}\\
  \cd{v\_sc\_electrico\_corregido}\\ \cd{v\_sc\_radiacion\_corregida}\\
  \nota{85\,°C, rangos, offset, pre-2025}};

\node[comp,text width=46mm,right=14mm of corr] (cal) {\textbf{Capa 2 --- calibración}\\
  \cd{v\_sc\_radiacion\_calibrada}\\
  \nota{W/m$^2$, kt*, \cd{qc\_ok}}};

\node[comp,text width=46mm,below=12mm of cal] (perf) {\textbf{Capa 3 --- desempeño}\\
  \cd{v\_sc\_performance}\\
  \nota{\cd{pr\_pv1}, \cd{pr\_pv2}}};

\node[comp,text width=42mm,below=12mm of corr] (cs) {\cd{radiacion\_sc\_clearsky}\\ \cd{radiacion\_sc\_poa}\\
  \nota{calculadas con pvlib}};

\draw[flecha] (crudo) -- node[etiq,above]{vista} (corr);
\draw[flecha] (corr) -- node[etiq,above]{vista} (cal);
\draw[flecha] (cs) -- (cal);
\draw[flecha] (cs) -- (perf);
\draw[flecha] (corr.south) |- (perf.west);

\node[etiq,below=3mm of crudo,align=center,fill=none]
  {sin pérdida de información:\\ nada se sobrescribe};

\end{tikzpicture}
}
\caption{Capas del modelo. Una consulta puede pedir el dato tal como salió del sensor o
el dato corregido y calibrado, según lo que necesite justificar.}
\end{figure}

\subsection{Relaciones}

\begin{center}
\small
\begin{tabular}{@{}p{44mm}p{18mm}p{74mm}@{}}
\toprule
\textbf{Relación} & \textbf{Tipo} & \textbf{Contenido} \\
\midrule
\cd{monitoreo\_sc\_electrico} & tabla & Crudo eléctrico a 5\,min: voltaje, corriente,
  potencia y energía por arreglo, salida AC, temperatura del inversor y de cada arreglo. \\
\cd{radiacion\_sc\_15s} & tabla & Crudo de radiación a 15\,s: incidente, reflejada,
  albedo, y las mismas del piranómetro SP722. \\
\cd{radiacion\_sc\_clearsky} & tabla & \cd{cs\_ghi\_wm2} por instante, calculado con pvlib. \\
\cd{radiacion\_sc\_poa} & tabla & Irradiancia en el plano del arreglo, frontal y
  bifacial, para PV1 y PV2. \\
\cd{diccionario\_variables} & tabla & Definición en prosa de cada variable. \\
\cd{\_ingest\_log} & tabla & md5 y conteo por archivo procesado. \\
\midrule
\cd{v\_sc\_electrico\_corregido} & vista & Aplica rangos válidos; fuera de rango pasa a NULL. \\
\cd{v\_sc\_radiacion\_corregida} & vista & Neutraliza el \emph{offset}, recorta negativos
  y anula el período inválido; marca \cd{valido}. \\
\cd{v\_sc\_radiacion\_calibrada} & vista & Radiación en W/m$^2$, \cd{kt\_star} y
  bandera \cd{qc\_ok}. \\
\cd{v\_sc\_performance} & vista & \cd{pr\_pv1} y \cd{pr\_pv2}: Performance Ratio por arreglo. \\
\midrule
\cd{lecturas\_ambientales\_sc} & tabla & Store de la ingesta ambiental, formato largo. \\
\cd{predicciones} & tabla & Auditoría de cada pronóstico emitido. \\
\cd{agente\_log} & tabla & Log estructurado del ETL y de los agentes. \\
\cd{gasto\_diario} & tabla & Gasto acumulado en el LLM, por día. \\
\bottomrule
\end{tabular}
\end{center}

\subsection{Reglas de corrección}

Los umbrales son parámetros en \cd{config.py} y el SQL de las vistas se genera a partir
de ellos, así que cambiar un límite es cambiar una constante y regenerar el DDL.

\begin{center}
\small
\begin{tabular}{@{}p{40mm}p{34mm}p{62mm}@{}}
\toprule
\textbf{Variable} & \textbf{Rango válido} & \textbf{Motivo} \\
\midrule
Temperaturas & 10--80\,°C & Criterio de Leo Cardinale; además el valor exacto 85{,}0
  es el código de error del DS18B20 desconectado. \\
Potencia por string & 0--5000\,W & Cota física provisional; el instalado real es
  1420\,Wp por arreglo. \\
Voltaje de string & 0--600\,V & Rango del inversor. \\
Corriente de string & 0--20\,A & Rango del inversor. \\
Frecuencia de red & 55--65\,Hz & Red de 60\,Hz. \\
Tensión AC & 100--280\,V & Salida monofásica. \\
Albedo & 0--1 & Definición del cociente. \\
Irradiancia & no aplica & \emph{Offset} $-38{,}845 \to 0$; negativos $\to 0$; todo lo anterior al
  2025-07-01 pasa a NULL (error de sensor corregido a mediados de 2025). \\
\bottomrule
\end{tabular}
\end{center}

\subsection{Calibración y desempeño}

No hay constante de calibración de fábrica: lo que se creía una ``celda calibrada'' resultó
ser un nombre comercial. Se calibra entonces contra el cielo despejado, comparando lo que
mide el sensor con lo que debería medir un día sin nubes. Los datos del período válido ya
venían en W/m$^2$, así que la escala quedó en 1{,}0 y lo que aporta esta capa es el índice
de claridad y el control de calidad.

Toda la capa se reduce a dos fórmulas. La primera dice cuánta de la luz posible llegó
de verdad:

\begin{equation*}
  \mathrm{kt^*} = \frac{\mathrm{GHI_{medida}}}{\mathrm{GHI_{cielo\ despejado}}}
  \qquad\text{sólo donde}\quad \mathrm{GHI_{cs}} > 20\ \mathrm{W/m^2}
\end{equation*}

La segunda compara la potencia real con la que cabría esperar de la luz que recibe el
plano del arreglo:

\begin{equation*}
  \mathrm{PR} = \frac{P/P_0}{\mathrm{POA}/1000}
  \qquad P_0 = 1420\ \mathrm{Wp\ por\ arreglo}
\end{equation*}

Como los paneles son bifaciales, la POA suma las dos caras: frontal $+\ \varphi\cdot$trasera,
con $\varphi = 0{,}80$ y el albedo medido.

El resultado da confianza porque converge: dos arreglos con geometrías muy distintas
aterrizan casi en el mismo número, PV1 inclinado en $0{,}622$ y PV2 vertical en $0{,}626$.

\begin{clave}[title=Geometría confirmada del sitio]
1420\,Wp por arreglo ($4 \times 355$\,Wp), 2840\,Wp en total, ambos bifaciales.
PV1 es el arreglo inclinado (\emph{tilt} 20°, azimut 150°); PV2 es el vertical
(\emph{tilt} 90°, azimut 50°), con Norte $=0$° y sentido horario positivo. Zona horaria
UTC$-6$ sin horario de verano.
\end{clave}
